\documentclass[11pt]{article}
\usepackage{amsmath,amssymb,amsthm}
\usepackage{geometry}
\usepackage{hyperref}
\usepackage{microtype}
\usepackage{booktabs}
\geometry{margin=1.2in}

\newtheorem{theorem}{Theorem}[section]
\newtheorem{proposition}[theorem]{Proposition}
\newtheorem{definition}[theorem]{Definition}
\newtheorem{remark}[theorem]{Remark}

\newcommand{\K}{\mathrm{K}}
\newcommand{\MDL}{\mathrm{MDL}}
\newcommand{\R}{\mathbb{R}}

\title{Science Has an Objective Function}
\author{Greg Coppola}
\date{2026}

\begin{document}
\maketitle

\begin{abstract}
Science has always had an objective function. Written down, it is:
minimize $L(H) + L(D \mid H)$, where $H$ is a hypothesis, $D$ is
observed data, $L(H)$ is the description length of the hypothesis
(approximating its Kolmogorov complexity), and $L(D \mid H)$ is the
description length of the data given the hypothesis (the negative
log-likelihood). This is the Minimum Description Length (MDL)
principle---the computable approximation to Solomonoff induction.
We make four contributions. First, we trace the history of philosophy
of science as a single argument converging on this framework, and
identify ten cases where independent traditions discovered the same
structure in different language. Second, we argue that the Quantified
Boolean Bayesian Network (QBBN) is the best known candidate for a
universal practical coding scheme for the MDL objective: it is
complete over first-order logic, generative, tractable via belief
propagation, and---by the \texttt{shannon} result---the coding scheme
that transformers implicitly implement. Third, we demonstrate the
framework empirically in two domains: programs and code
(\texttt{auto-research}), where a new result shows current AutoResearch
systems are one $L(H)$ term short of the full objective, accounting
for a 65.2\% improvement on benchmark problems; and logical
proposition datasets (\texttt{mdl-science}), where QBBN-coded MDL
selects correct conjunctive structure, amortizes 8 bits of theory over
510 bits of savings at scale, and exhibits the discrete regularity
crossover that instantiates Kuhn's paradigm shift at the bit level.
Fourth, we argue that a system with the right objective function is
doing science regardless of substrate: automated theorem proving with
MDL is not adjacent to AGI---it is AGI. The objective function is MDL.
We already know what it is.
\end{abstract}

\tableofcontents
\newpage

\section{Introduction}
\label{sec:introduction}

Science has always had an objective function. It has never been written
down in a single place. This paper writes it down.

The objective function is:
\[
  \text{minimize} \quad L(H) + L(D \mid H)
\]
where $H$ is a hypothesis, $D$ is observed data, $L(H)$ is the
description length of the hypothesis, and $L(D \mid H)$ is the
description length of the data given the hypothesis. This is the
Minimum Description Length (MDL) principle---the computable
approximation to Solomonoff induction, which is Bayesian inference
with the Kolmogorov/Solomonoff prior. It is not a new proposal. It is
a formal statement of what science has always been doing informally.

This paper makes four contributions.

\textbf{A synthesis of the history of philosophy of science.} The
history of philosophy of science from Hume through Popper, Kuhn,
Lakatos, Feyerabend, and the Bayesian tradition is a single long
argument about one problem: Hume's problem of induction. Every major
movement is a partial approximation to the framework above. We trace
this arc and show that each movement captured one component of the
full solution without being able to state it. Section~\ref{sec:history}
develops this argument. Section~\ref{sec:connections} identifies ten
cases where ideas from independent traditions---philosophy,
computability theory, information theory, neuroscience, machine
learning---turn out to be the same idea expressed in different language.

\textbf{The QBBN as a universal coding scheme for logical hypotheses.}
MDL requires a coding scheme: a way of assigning description lengths to
hypotheses. The Kolmogorov prior gives the theoretically optimal
scheme, but it is uncomputable. Every practical coding scheme covers
only a restricted hypothesis class. The Quantified Boolean Bayesian
Network (QBBN)~\cite{coppola2024qbbn} is the best known candidate for
a universal practical coding scheme: it is complete over first-order
logic (covering all of mathematics and science), generative (assigning
likelihoods to arbitrary proposition datasets), and tractable (belief
propagation runs in polynomial time). Moreover, the
\texttt{shannon}~paper~\cite{coppola2026} proves that transformers
implement exact belief propagation on QBBN-structured factor graphs,
connecting the coding scheme to the architecture that currently
dominates machine learning. Section~\ref{sec:coding} develops this
argument.

\textbf{Two empirical demonstrations.} The framework is demonstrated
empirically in two domains. In the domain of code and programs
(Section~\ref{sec:experiments-code}), we demonstrate compression
ratios, MDL model selection, the replication crisis as a missing
$L(H)$ term, the search/verify asymmetry, and the Karpathy
gap---current AutoResearch systems minimize $L(D \mid H)$ without
$L(H)$, and the missing term accounts for a 65.2\% improvement on
benchmark problems. In the domain of logical propositions
(Section~\ref{sec:experiments-logic}), we use the QBBN coding scheme
to show that conjunctive structure wins over disjunctive, that
theories are amortized over data (8 bits of theory saves 510 bits at
$n = 1000$), and that MDL selects correct theory complexity as a
function of data regularity, with a discrete crossover that
instantiates Kuhn's paradigm shift mechanism at the bit level.

\textbf{Theorem proving is science; the objective function is AGI.}
The critics who say ``automated theorem proving is not AGI because it
is not science'' are relying on an implicit theory of science that has
no formal content. Under the correct theory, the verification oracle
is interchangeable: a particle accelerator, a type-checker, and an
MDL score calculation are all instances of the same abstract
structure. A system with the right objective function is doing science
regardless of substrate. Section~\ref{sec:agi} develops this argument
and proposes a formal definition of AGI grounded in the framework.

\subsection*{Related Work}

The Bayesian/Kolmogorov framework synthesized here draws on
Solomonoff's universal prior~\cite{solomonoff1964}, Kolmogorov
complexity~\cite{kolmogorov1965}, Rissanen's MDL
principle~\cite{rissanen1978}, and the Bayesian tradition in
philosophy of science~\cite{howson2006}. The connection between
paradigm shifts and MDL phase transitions extends
Kuhn~\cite{kuhn1962} and Lakatos~\cite{lakatos1978}. The
Bayesian brain and predictive coding framework draws on Rao and
Ballard~\cite{rao1999} and Friston~\cite{friston2005}. The QBBN
coding scheme is developed in~\cite{coppola2024qbbn}. The proof that
transformers implement exact belief propagation on QBBN-structured
factor graphs is in~\cite{coppola2026}. The AutoResearch system
discussed in the Karpathy gap analysis is~\cite{karpathy2026}.

The \texttt{auto-research} and \texttt{mdl-science} repositories
contain all code and data for the empirical demonstrations. All
results are reproducible by following the instructions in those
repositories.
\section{The Objective Function}

\subsection{The Claim}

Science finds short descriptions of large bodies of observations. The right
way to formalize ``short description'' is Kolmogorov complexity. The right
way to formalize ``finds'' is Bayesian inference with the Solomonoff prior.
Everything else in this paper is unpacking those three sentences.

The formal statement. Let $H$ be a hypothesis---a theory, a model, a
program---and let $D$ be observed data. The objective function for science is:
\begin{equation}
  \text{maximize} \quad \log P(H \mid D)
  = \log P(D \mid H) + \log P(H) + \text{const}
  \label{eq:bayes}
\end{equation}
by Bayes' rule, where $P(D \mid H)$ is the likelihood and $P(H)$ is the
prior. This is uncontroversial as a framework. The only question is: what
prior?

\subsection{The Prior Problem}

If the prior $P(H)$ can be chosen arbitrarily, Bayes' rule is a framework,
not a complete theory. Two scientists with different priors will reach
different posteriors from the same data. The standard objection to
Bayesianism is that this makes science subjective.

The Kolmogorov/Solomonoff answer: the prior is not arbitrary. There is a
unique, principled choice grounded in the theory of computation.

\subsection{The Solomonoff Prior}

The Kolmogorov complexity $\K(x)$ of a string $x$ is the length of the
shortest program (on a fixed universal Turing machine $U$) that outputs $x$
and halts:
\[
  \K(x) = \min\{ |p| : U(p) = x \}.
\]
Solomonoff~\cite{solomonoff1964} proposed assigning to each hypothesis $H$
the prior:
\[
  P(H) \propto 2^{-\K(H)}.
\]
Simpler hypotheses---those with shorter generators---receive exponentially
higher prior probability. This is Occam's Razor, derived from first
principles rather than assumed as a philosophical preference.

The Solomonoff prior is \emph{universal}: for any computable probability
distribution $P$ over strings, there exists a constant $c_P$ such that
$M(x) \geq c_P \cdot P(x)$ for all strings $x$, where $M$ is the Solomonoff
measure. No computable prior assigns more probability to any hypothesis than
$M$ does, up to a fixed factor. The prior dominates every computable
alternative.

\subsection{Minimum Description Length}

Because $\K$ is uncomputable (Section~\ref{sec:kolmogorov}), the practical
version is the Minimum Description Length (MDL) principle~\cite{rissanen1978}.
Instead of the shortest program, use the shortest description under a fixed
code. The objective becomes:
\begin{equation}
  \text{minimize} \quad L(H) + L(D \mid H)
  \label{eq:mdl}
\end{equation}
where $L(H)$ is the description length of the hypothesis and $L(D \mid H)$
is the description length of the data given the hypothesis (the negative
log-likelihood under a specific coding scheme).

The two terms trade off directly. $L(H)$ penalizes complexity: a theory that
is too elaborate spends too many bits describing itself. $L(D \mid H)$
penalizes poor fit: a theory that predicts the data badly leaves many residual
bits to transmit. The optimum is the theory that compresses the combined
description most. Science is the search for short two-part codes for the world.

\subsection{Occam's Razor, Derived}

A consequence worth stating explicitly: two theories that predict the data
equally well are not equally good. The simpler one has higher posterior. A
more complex theory must predict the data \emph{better} to compensate for its
higher complexity penalty. The scientist who adds epicycles to save a theory is
doing something formally wrong: increasing $\K(H)$ without a commensurate gain
in $\log P(D \mid H)$.

The objective function is \emph{objective} in both senses of the word: it is
a target (something to minimize) and it is observer-independent (grounded in
the structure of computation, not in any particular scientist's preferences).
The Solomonoff prior is invariant across observers up to a fixed constant
determined by the choice of universal Turing machine---not by whoever is doing
the science.
\section{Kolmogorov Complexity}
\label{sec:kolmogorov}

\subsection{The Intuition}

Consider two strings of a thousand bits. String A alternates 0 and 1
indefinitely; String B is a typical output of a fair coin. Both are a
thousand bits long. String A can be described in a few dozen characters:
``alternating 0 and 1, one thousand times.'' String B, if genuinely random,
cannot be described more briefly than writing it out in full.

Kolmogorov complexity makes this precise. The complexity $\K(x)$ of a string
$x$ is the length of the shortest computer program that outputs $x$ and halts.
String A has very low $\K$; String B has $\K$ close to its own length.

Three examples locate the concept concretely. The digits of $\pi$ look random
but have very low $\K$: any standard algorithm for computing $\pi$ provides a
program of a few hundred bytes that generates arbitrarily many digits.
Newton's law of universal gravitation $F = Gm_1 m_2 / r^2$ is a program of a
few dozen bytes that predicts the trajectories of every planet, moon, and comet
we have observed, to many decimal places. A random genome, by contrast, would
have $\K$ close to its raw length---but the actual human genome is substantially
more compressible, because evolution is a compressive process.

\subsection{The Formal Definition}

Fix a universal Turing machine $U$. The Kolmogorov complexity of a finite
binary string $x$ with respect to $U$ is:
\[
  \K_U(x) = \min\{ |p| : U(p) = x \}
\]
where $|p|$ is the length of program $p$ in bits. In practice, the
prefix-free variant (where no valid program is a prefix of another) is
standard; it has cleaner information-theoretic properties.

The key theorems:

\begin{theorem}[Invariance~\cite{kolmogorov1965}]
For any two universal Turing machines $U$ and $V$:
\[
  |\K_U(x) - \K_V(x)| \leq c_{U,V}
\]
where $c_{U,V}$ depends only on the machines, not on $x$. For large strings,
this constant is negligible relative to $\K(x)$.
\end{theorem}

The Invariance Theorem is what makes $\K(x)$ objective: the choice of machine
changes the value by at most a fixed overhead, independent of the string being
measured. For any two computable descriptions, the complexity is the same up
to a translator constant.

\begin{proposition}
For any length $n$, at most $2^{n-c}$ strings of length $n$ have
$\K(x) < n - c$. Most strings are incompressible.
\end{proposition}

This is a counting argument: there are only $2^{n-c}$ programs shorter than
$n-c$ bits. When we encounter a compressible string in nature---a genome, a
physical trajectory, a spectrum---that compressibility is evidence of an
underlying generative process. Structure is rare.

\subsection{Uncomputability: A Feature, Not a Bug}

$\K(x)$ is not computable. No program, given $x$ as input, reliably outputs
$\K(x)$ and halts. The proof is a formalization of Berry's paradox: supposing
$\K$ were computable, one could construct a short program that outputs a string
with high complexity---a contradiction.

This is sometimes presented as a defect. It is not. Kolmogorov complexity is
the \emph{right} measure of structural complexity. Its uncomputability is the
signature of universality: $\K$ captures structure defined by all possible
programs, not any particular one. If $\K$ were computable, it would be
reducible to some specific algorithm, and that algorithm would make implicit
choices about what structure counts.

The scientist's relationship to $\K$ is the same as the mathematician's
relationship to real numbers: most are uncomputable, yet the reals are the
right framework. Scientists approximate $\K$ from above. Every proposed
theory is an upper bound on $\K(D)$. The search for shorter descriptions is
the search for a tighter bound. The fact that the true minimum is uncomputable
is precisely why the search never ends.

\subsection{MDL as the Computable Approximation}

The Minimum Description Length principle~\cite{rissanen1978} replaces $\K$
with a description length under a fixed coding scheme. For a parametric model
$H$ with parameters $\theta$:
\[
  L(H) + L(D \mid H) = L(\theta) - \log P(D \mid \theta)
\]
This is computable for any fixed model class and coding scheme. MDL is
model-selection consistent: as data grows, it converges to the true model
when the true model is in the hypothesis class. It reduces overfitting by
penalizing complexity directly, without ad hoc regularization.

The relationship to Bayes is exact. Under a prior $P(H) \propto 2^{-L(H)}$:
\[
  \log P(H \mid D) = -L(H) + \log P(D \mid H) + \text{const} = -\MDL(H,D) + \text{const}
\]
Maximizing the Bayesian posterior is equivalent to minimizing MDL, for the
right choice of prior and coding scheme.

\subsection{Connections to Information Theory}

The expected Kolmogorov complexity of a draw from a computable distribution
$P$ equals its Shannon entropy $H(P)$ up to lower-order terms:
\[
  \mathbb{E}[K(X)] \approx H(X).
\]
Shannon entropy is Kolmogorov complexity averaged over a distribution.
Kolmogorov complexity is the pointwise version: it applies to individual
strings, without requiring a distribution. When you have a single string and
want to know how much structure it has---without committing to a probability
model---Kolmogorov complexity is the right tool.

Lossless compression (gzip, bzip2, zstd) provides empirical upper bounds on
$\K$. A string that compresses well has low $\K$; one that does not compress
has high $\K$. This means compression ratio is a practical proxy for
structural complexity, usable without specifying any explicit model.
\section{The Coding Scheme Problem}
\label{sec:coding}

\subsection{MDL Requires a Coding Scheme}

The MDL objective is:
\[
  \text{minimize} \quad L(H) + L(D \mid H)
\]
This formula is only well-defined once you specify how hypotheses and
data are encoded---what ``description length'' means. Different coding
schemes give different numerical values of $L(H)$ and $L(D \mid H)$,
and in principle different rankings of theories.

The Kolmogorov framework provides the theoretically correct answer:
the description length of a hypothesis is $\K(H)$, the length of the
shortest program that generates it. The Solomonoff prior
$P(H) \propto 2^{-\K(H)}$ is the resulting prior. This coding scheme
is universal: it dominates every computable distribution up to a
multiplicative constant (Section~\ref{sec:kolmogorov}). For large
enough data, any two universal coding schemes give the same ranking
of hypotheses---the constant washes out.

The problem is that $\K(H)$ is uncomputable. You cannot implement the
Kolmogorov coding scheme directly. Every practical MDL application
uses a restricted coding scheme: a parametric family, a grammar, a
neural architecture. These are computable but not universal---they
cover a restricted hypothesis class and give no principled way to
compare hypotheses from different families.

This is the coding scheme problem: to have a complete formal theory of
science, we need a single coding scheme that is simultaneously
\emph{universal} (covering all scientific hypotheses),
\emph{generative} (assigning likelihoods to data), and
\emph{tractable} (computable in practice).

\subsection{What a Universal Practical Coding Scheme Requires}

A coding scheme for scientific hypotheses must satisfy four
properties to serve as a universal practical approximation to the
Kolmogorov prior:

\begin{enumerate}

\item \textbf{Universality.} The hypothesis space must cover all
scientifically expressible claims. Since all of mathematics and
science can be expressed in first-order logic
(FOL)~\cite{coppola2024qbbn}, a coding scheme is universal in the
relevant sense if it can represent any FOL formula.

\item \textbf{Generativity.} The coding scheme must define a joint
probability distribution over propositions---it must assign
$P(D \mid H)$ for any dataset $D$ and hypothesis $H$. Without
generativity, you cannot compute $L(D \mid H)$ and the MDL objective
is not defined. This rules out pure theorem provers, which give truth
values but not probabilities.

\item \textbf{Tractability.} The likelihoods $P(D \mid H)$ must be
computable in reasonable time. Exact inference in general Bayesian
networks is \#P-complete; a universal scheme needs a tractable
approximation.

\item \textbf{Structural transparency.} The coding scheme should make
the structure of a hypothesis explicit---which parts are the hard
logical constraints and which parts are the learned statistical
weights. This is what allows $L(H)$ to be decomposed into a structure
cost and a weight cost, connecting the coding scheme to the
Kolmogorov prior's notion of program length.

\end{enumerate}

No existing system satisfies all four. Pure FOL theorem provers
satisfy (1) and (4) but not (2) or (3). Markov Logic
Networks~\cite{coppola2024qbbn} satisfy (1) and (2) but not (3):
inference is \#P-complete and requires MCMC approximation. Neural
networks satisfy (2) and (3) but not (1) or (4): they approximate
logical structure without representing it explicitly, and their
hypothesis space is the space of weight configurations, not logical
propositions.

\subsection{The QBBN as the Best Known Candidate}

The Quantified Boolean Bayesian Network
(QBBN)~\cite{coppola2024qbbn} is the best known candidate for a
universal practical coding scheme. It satisfies all four properties.

\textbf{Universality (FOL completeness).} The QBBN paper proves that
any first-order logical formula can be expressed as a QBBN factor
graph. The QBBN is complete over FOL in the same sense that FOL is
complete over mathematics and science: any claim that can be stated
precisely can be stated as a QBBN.

\textbf{Generativity.} The QBBN defines a joint probability
distribution over boolean propositions via its factor graph. The AND
gates are deterministic (they implement exact boolean conjunction);
the OR gates are learned linear exponential models that assign
conditional probabilities. Together they define $P(D \mid H)$ for any
dataset $D$ of boolean observations.

\textbf{Tractability.} Inference in the QBBN uses belief propagation
(BP) on the AND/OR bipartite factor graph. Exact BP is polynomial on
trees; on loopy graphs it is an approximation, but the QBBN paper
reports empirical convergence in all tested cases. The bipartite
AND/OR structure---alternating conjunction and disjunction
layers---is the key architectural choice that makes BP tractable:
conjunction nodes are deterministic (no learned parameters, $O(2^n)$
but constant in practice for small conjuncts), and disjunction nodes
admit $O(n)$ updates under the Noisy-OR approximation.

\textbf{Structural transparency.} The AND/OR decomposition makes the
structure of a hypothesis explicit. AND gates represent the hard
logical constraints---the ``program'' part of the hypothesis. OR gates
represent the learned statistical weights---the ``parameter'' part.
The description length of the hypothesis is:
\[
  L(H) = L(\text{AND structure}) + L(\text{OR weights})
\]
The AND structure is the set of implication links---which premises
conjoin to produce which conclusions. Each link costs a fixed number
of bits to specify (its type and participating propositions). The OR
weights are the conditional probabilities learned for each disjunction
node; each weight costs $w$ bits to encode at $w$-bit precision. This
maps directly onto the Kolmogorov prior: AND structure corresponds to
program structure (fixed overhead per instruction), OR weights
correspond to program parameters (variable, proportional to
precision).

The \texttt{mdl-science} experiments in
Section~\ref{sec:experiments-logic} instantiate this coding scheme
directly. Each experiment uses QBBN theories---implication links and
conjunctive links---and computes MDL scores that reflect the true
$L(H) + L(D \mid H)$ under QBBN coding. The results (conjunctive
structure wins, theories are amortized, MDL selects correct
complexity) are consequences of the coding scheme's alignment with the
true generative structure.

\subsection{The Transformer as Independent Confirmation}

The coding scheme argument gains additional force from an independent
direction: the \texttt{shannon} paper~\cite{coppola2026} proves that
transformers implement exact belief propagation on QBBN-structured
factor graphs.

Specifically, the paper shows that:
\begin{itemize}
  \item Attention heads implement AND-type joint beliefs over
        query-key pairs---the AND gate computation of the QBBN.
  \item Feed-forward networks implement OR-type disjunctive
        updates---the OR gate computation of the QBBN.
  \item The residual stream is the shared message-passing bus---the
        variable nodes of the QBBN factor graph.
\end{itemize}

This means that a transformer trained by gradient descent on
next-token prediction is implicitly minimizing $L(D \mid H)$ where
$H$ is parameterized as a QBBN factor graph. The QBBN coding scheme
is not a design choice imposed on top of transformer training---it is
what transformer training discovers.

This has a direct implication for the Karpathy gap
(Section~\ref{sec:experiments-code}). The gap we demonstrated on toy
sort programs---full MDL achieves 65.2\% improvement while
$L(D \mid H)$-only achieves 0\%---is not a toy result. It applies
to every transformer trained without an explicit $L(H)$ penalty. A
language model trained to minimize validation loss (= $L(D \mid H)$
in QBBN coding) without any constraint on $L(H)$ will accumulate
complexity in the OR weights whenever doing so does not hurt
validation loss. The weight matrix grows without any signal that
something is wrong. Adding an explicit $L(H)$ penalty---compressed
size of the weight matrix, added to the training objective---would
complete the MDL objective in QBBN coding.

\subsection{The Open Problem}

The QBBN satisfies the four properties above for the domain of
logical/propositional hypotheses. It does not solve the general
coding scheme problem.

The open problem is cross-domain comparison. The QBBN gives a
principled MDL ranking of QBBN theories against each other. It does
not give a principled MDL ranking of a QBBN theory against, say, a
differential equation or a statistical regression model---because
there is no shared coding scheme that covers all three. Such a ranking
exists in principle: the Kolmogorov prior assigns description lengths
to all computable hypotheses, and the ranking is well-defined up to a
constant. But computing it requires $\K(H)$, which is uncomputable.

This is the computable residue of Hume's problem
(Section~\ref{sec:connections}). Hume showed that induction cannot be
deductively justified; we showed (in the Hume/Halting connection) that
this is equivalent to the uncomputability of $\K$. The coding scheme
problem is the same impossibility in engineering form: you cannot
build a universal computable coding scheme for all scientific
hypotheses for the same reason you cannot solve the halting problem.

What you can do---and what this paper does---is:
\begin{enumerate}
  \item Identify the best known candidate for a universal coding
        scheme within the logical/propositional domain (QBBN).
  \item Show that this coding scheme connects to the dominant
        practical architecture (transformers) via the
        \texttt{shannon} result.
  \item Demonstrate empirically that MDL in QBBN coding selects
        correct structure, penalizes spurious complexity, and
        amortizes over data as the theory predicts.
  \item Name the open problem honestly: a universal computable coding
        scheme covering all scientific domains does not yet exist.
\end{enumerate}

The Kolmogorov framework tells us such a scheme would exist if $\K$
were computable. It also tells us $\K$ is not computable. The QBBN is
the closest computable approximation currently known for the domain
where science and logic intersect---which, given FOL's expressive
completeness over mathematics and science, is most of the domain that
matters.
\section{The History of Philosophy of Science as a Single Argument}

The history of philosophy of science is a single long argument about one
problem: Hume's problem of induction. Every major movement is either a
response to it, a refinement of a previous response, or an attempt to
dissolve it. The Bayesian/Kolmogorov framework is where the argument arrives.
We trace the through-line here, briefly; each movement clarifies what the
framework is solving.

\subsection{Hume: The Impossibility}

David Hume~\cite{hume1748} identified the deepest problem in the philosophy
of science. The problem of induction: no finite set of observations logically
justifies a universal law. The inference from ``all observed $F$s are $G$s''
to ``all $F$s are $G$s'' is not deductively valid. You can observe a million
white swans; the next one might be black. It was, in Australia.

Worse: the obvious response---induction is justified because it has worked in
the past---is itself an inductive argument. The justification is circular.
Hume concluded that inductive inference is a matter of psychological habit,
not reason.

\subsection{Popper: The Right Asymmetry, the Wrong Conclusion}

Karl Popper~\cite{popper1934} noticed that confirmation is logically weak
while falsification is logically strong. No number of white swans proves the
law; one black swan refutes it. He proposed falsifiability as the
demarcation criterion for science and characterized scientific progress as
conjecture and refutation: bold hypotheses, then honest attempts to falsify
them.

Popper was right about the asymmetry and right that unfalsifiable theories
are scientifically empty. He was wrong to conclude that science involves no
induction. His ``corroboration'' measure---how well a theory has survived
tests---is, precisely defined, the Bayesian likelihood ratio $P(E \mid H) /
P(E \mid \neg H)$: the update factor, stripped of the prior he refused to use.
Popper computed the right quantity and declined to multiply by the prior.
The Solomonoff prior completes his programme.

\subsection{Kuhn: The Right Phenomenology, No Normative Theory}

Thomas Kuhn~\cite{kuhn1962} argued that Popper's account does not match how
science actually works. Scientists operate within \emph{paradigms}---shared
frameworks of assumptions and methods---and do not abandon them at the first
anomaly. When anomalies accumulate to the point where confidence breaks down,
a revolution occurs: one paradigm replaces another, discontinuously.

Kuhn correctly described the phenomenology. He lacked a normative account:
when \emph{should} a paradigm shift occur? The MDL framework supplies the
answer (Section~\ref{sec:connections}).

\subsection{Lakatos: The Right Criterion, Incompletely Formalized}

Imre Lakatos~\cite{lakatos1978} synthesized Popper and Kuhn via the
methodology of scientific research programmes. Each programme has a
\emph{hard core} of protected assumptions and a \emph{protective belt} of
auxiliary hypotheses that absorb anomalies. A programme is \emph{progressive}
if its protective belt modifications generate new confirmed predictions;
\emph{degenerating} if they only explain known anomalies without predicting
new ones.

Under MDL, a progressive programme is one where successive modifications
decrease $L(H) + L(D \mid H)$ on new data. A degenerating programme is one
where modifications increase $L(H)$ without decreasing $L(D \mid H)$ on new
observations. Lakatos gave the right criterion; MDL makes it computable.

\subsection{Feyerabend: Methodological Anarchism as the Solomonoff Prior}

Paul Feyerabend~\cite{feyerabend1975} argued that no hypothesis should be
ruled out \emph{a priori}: science has always advanced by entertaining ideas
that violated current methodological rules. ``Anything goes'' as a prior
policy.

The Solomonoff prior assigns nonzero probability to every computable
hypothesis. Nothing is excluded before the data speaks. Feyerabend's
permissiveness is not anarchism---it is the correct prior policy. A prior
that excludes any computable hypothesis is provably suboptimal. What
Feyerabend called anarchism, Solomonoff called universality.

\subsection{Bayesian Confirmation: The Culmination}

Bayesian confirmation theory~\cite{howson2006} provides the most complete
response to Hume. Evidence $E$ confirms $H$ if and only if $P(H \mid E) >
P(H)$; the degree of confirmation is the likelihood ratio $P(E \mid H) /
P(E)$. This dissolves the raven paradox (observing a red fire engine confirms
``all ravens are black,'' but only infinitesimally), handles old evidence via
counterfactual reasoning, and validates Lakatos's distinction as a distinction
between increasing and decreasing marginal likelihood.

Carnap spent his career trying to construct an inductive logic---a formal
system that would determine the degree of confirmation of any hypothesis by
any evidence. He never succeeded because he had no principled account of the
prior. The Solomonoff prior completes the programme he started.
\section{Ten Identifications}
\label{sec:connections}

The directories of the \texttt{science} repo were built separately. This
section is about what happens when those threads cross. The cases below are
not analogies---loose, suggestive comparisons. They are \emph{identifications}:
cases where two ideas from different traditions turn out to be the same idea,
expressed in different language, developed by people who were not talking to
each other. The convergence is evidence that something real is being pointed at.

\subsection{Hume's Problem Is the Halting Problem}

Hume showed that induction cannot be logically justified: you cannot know,
from finite observations, that a universal law holds. Turing showed that
you cannot decide, in general, whether a program halts.

These are the same result. A universal law is a short program---a generator
for an infinite class of observations. Confirming the law means certifying
that the proposed program is the shortest consistent with the data. But
certifying that no shorter program exists requires checking whether every
shorter program eventually fails to match the data---instances of the halting
problem. Since the halting problem is undecidable, $\K(x)$ is uncomputable,
and induction cannot be deductively justified.

The reason you cannot justify induction is the same reason $\K(x)$ is
uncomputable. Hume's problem is not a philosophical puzzle awaiting a
philosophical solution. It is an undecidability result, and it has the same
status as the halting problem: provably unsolvable, not merely unsolved.

What can be done---and what the Bayesian/Kolmogorov framework does---is show
that the impossibility of certainty does not prevent optimal inference.
Solomonoff induction is provably optimal for prediction in any computable
environment. The gap between ``optimal'' and ``certain'' is exactly the gap
between computable approximations to $\K(x)$ and $\K(x)$ itself.

\subsection{Paradigm Shifts Are MDL Phase Transitions}

Kuhn described paradigm shifts as discontinuous and partly irrational. Under
MDL, they are mathematically inevitable and perfectly rational.

A paradigm is a hypothesis class---a family of theories sharing common
structure. Normal science is optimization within the class: adjusting
parameters, adding auxiliary hypotheses, refining the model. Each adjustment
is legitimate as long as it decreases total description length.

A paradigm shift occurs when a competing hypothesis class achieves lower
$L(H) + L(D \mid H)$ than the incumbent with all its protective belt
modifications. The shift is discontinuous because hypothesis classes are
discrete objects: you cannot smoothly interpolate between Ptolemaic and
Copernican astronomy. But it is not irrational. The new paradigm wins because
it compresses the data better.

The discontinuity is a phase transition in the physics sense: a sharp change
in the optimal state of the system (which paradigm minimizes description
length) as a continuous parameter (accumulated data) crosses a threshold.
Kuhn correctly identified the phenomenology of phase transitions in
scientific knowledge. MDL provides the thermodynamics underneath.

\subsection{Popper's Corroboration Is Bayesian Updating Without the Prior}

Popper's corroboration measure, precisely defined~\cite{popper1954}, is:
\[
  C(H, E) = \frac{P(E \mid H) - P(E)}{P(E \mid H) - P(E \cdot H) + P(E)}
\]
This is a normalized likelihood ratio: how much more probable $E$ is under
$H$ than under the background probability of $E$. It is purely a function of
likelihoods. It does not involve the prior $P(H)$.

In Bayesian terms, Popper was computing exactly half of the update:
\[
  P(H \mid E) \propto \underbrace{P(E \mid H)}_{\text{Popper's corroboration}} \cdot \underbrace{P(H)}_{\text{refused}}
\]
He refused the prior because he found subjective probability unacceptable.
He was right to find it unacceptable; he was wrong to conclude the prior must
be abandoned. The Solomonoff prior $P(H) \propto 2^{-\K(H)}$ is objective:
a property of the hypothesis itself, invariant across observers. With it,
corroboration becomes posterior probability, and the refusal of induction
dissolves.

\subsection{The Replication Crisis Is an MDL Accounting Error}

The replication crisis in psychology and social science resulted from
p-hacking: testing many hypotheses and reporting only those achieving
$p < 0.05$. This is an MDL accounting error.

When a researcher tests $N$ hypotheses and reports the best, the true
description length of the reported hypothesis is not $L(H_{\text{reported}})$.
It is:
\[
  L(H_{\text{reported}}) + \log_2 N
\]
The extra $\log_2 N$ bits is the cost of the search: the information required
to specify which of the $N$ tested hypotheses was selected. Omitting this cost
produces apparent results that do not replicate, because the apparent
compression of the data was illusory.

The Bonferroni correction---adjusting the significance threshold by $N$---is
exactly the MDL accounting correction. It has always been the right thing to
do. The replication crisis is the empirical consequence of not doing it.

The reforms that followed---preregistration, larger samples, direct
replication, Bayesian model comparison---are all corrections toward honest
MDL accounting. Preregistration eliminates search cost from $L(H)$.
Preregistered studies have lower description length overhead, and their
results generalize because the $L(H)$ term is honest.

\subsection{Lakatos's Progressive vs.\ Degenerating Is Compression Growth vs.\ Decay}

Lakatos's distinction between progressive and degenerating research
programmes, translated into MDL:

A \emph{progressive} programme makes protective belt modifications that
decrease $L(H) + L(D \mid H)$ on new data. The modifications are short
(do not bloat description length) and genuinely compress new observations.

A \emph{degenerating} programme makes modifications that increase $L(H)$---
the hypothesis grows more complex---without decreasing $L(D \mid H)$ on new
observations. The programme adds complexity without adding compression.

Lakatos gave the right criterion for evaluating research programmes; MDL
makes it computable. The qualitative judgment (is this modification
progressive or defensive?) becomes a quantitative one (does
$\Delta L(H) + \Delta L(D \mid H)$ on new data decrease or increase?).

\subsection{Kant's Synthetic A Priori Is Lakatos's Hard Core}

Kant argued that some commitments---space, time, causality---are
preconditions of experience: they cannot be revised because they are what
makes experience possible. Lakatos argued that every research programme has
a hard core of commitments that are methodologically protected from
falsification.

These are the same idea. The difference is that Lakatos made them revisable
in principle: a research programme can be abandoned, hard core and all.
This is the correct update on Kant after non-Euclidean geometry falsified
his specific candidates (Euclidean space as a synthetic a priori). Lakatos
is Kant, empiricized.

Under the MDL framework, the hard core is the prior commitment to a
hypothesis class. It is protected not by necessity but by the cost of
switching to a different class (learning a new framework, retranslating
existing results). The protection is finite and can be overcome by sufficient
data. Kant's absolute preconditions become Lakatos's methodological
conventions, which become MDL's hypothesis-class priors.

\subsection{Feyerabend's ``Anything Goes'' Is the Solomonoff Prior}

Feyerabend argued that no methodological rule holds universally: every rule
that science has ever followed has been violated productively at some point.
``Anything goes'' as a description of scientific practice.

The Solomonoff prior assigns $P(H) \propto 2^{-\K(H)} > 0$ to every
computable hypothesis. No hypothesis is excluded before the data speaks.
This is precisely Feyerabend's permissiveness, formalized. A prior that
excludes any computable hypothesis---that declares some class of ideas
unscientific by methodological fiat---is provably suboptimal. The data
might support the excluded hypothesis, and a reasoner that cannot entertain
it incurs irreducible regret.

What Feyerabend called anarchism is, from the Kolmogorov perspective, the
universality of the prior. His insight was correct; his framing was
unnecessarily provocative.

\subsection{Hallucination and Perceptual Illusions Are the Same Failure Mode}

In predictive coding~\cite{rao1999}, the dominant framework in theoretical
neuroscience, the brain is a hierarchical generative system: higher areas
send predictions down the cortical hierarchy; lower areas send prediction
errors up. Perception is inference: the brain computes the posterior
distribution over world-states given sensory inputs, using priors built from
evolution and experience.

A perceptual illusion occurs when the prior is so strong that sensory
evidence cannot correct it. The hollow mask illusion is the canonical
example: a hollow face mask looks convex because the prior for convex faces
is overwhelming. The brain refuses to believe the correct interpretation.

In a language model, hallucination occurs when the prior over likely
continuations overwhelms the grounding signal from the context. The model
generates a continuation that is statistically plausible given training
distribution but inconsistent with the specific facts in the prompt.

Both are the same Bayesian failure: the prior dominates the likelihood. The
hollow mask illusion and language model confabulation are not analogous
phenomena. They are the same computational pathology, instantiated in
different substrates.

\subsection{The Bayesian Brain and the Transformer Are the Same Architecture}

Predictive coding models the brain as performing belief propagation on a
hierarchical factor graph~\cite{friston2005}: predictions propagate top-down,
prediction errors propagate bottom-up, and the system minimizes free energy
(variational approximation to surprise).

The transformer implements exact belief propagation on a specific factor
graph~\cite{coppola2026}: the attention mechanism computes AND-type joint
beliefs over query-key pairs; the feed-forward network computes OR-type
disjunctive updates; the residual stream is the shared message-passing bus.

These are not analogies. They are the same algorithm instantiated on
different graphs. The brain and the transformer are doing the same
computation. The differences are in the graph structure, the learning
algorithm, and the substrate---not in the fundamental operation.

\subsection{AutoResearch Is Solomonoff Induction Made Computable and Distributed}

Solomonoff induction is uncomputable but is the provably optimal prediction
strategy. Karpathy's AutoResearch system~\cite{karpathy2026} is a computable,
distributed approximation to it.

The correspondences are structural. The hypothesis space is the space of
training scripts and model configurations---a computable, enumerable space.
The reward signal (validation loss, benchmark accuracy) approximates
$P(D \mid H)$: lower loss means higher likelihood of the training data under
the modified hypothesis. The keep/discard decision approximates the Bayesian
posterior update. The agent's tendency to propose local modifications first
is an implicit Kolmogorov prior: smaller modifications correspond to lower
description length.

The key insight Karpathy identified---finding a good result is hard, but
verifying one is cheap---is exactly the search/verify asymmetry in Solomonoff
induction. The search over hypothesis space is expensive; checking whether a
proposed hypothesis compresses the data better is cheap. The distributed,
SETI@Home-style architecture parallelizes the search over hypothesis space.

Crucially, AutoResearch accepts contributions from untrusted sources. A
modification is evaluated by whether it reduces description length of the
training data---a property of the modification itself. The source is
irrelevant. This is the Kolmogorov principle applied to institutions:
\emph{trust the object, not the messenger}.
\section{Experiments I: Code and Programs}
\label{sec:experiments-code}

The experiments in this section instantiate the MDL objective function
in the domain of programs and code. All results are reproducible from
the \texttt{auto-research} repository. Six experiments, each isolating
one aspect of $L(H) + L(D \mid H)$: \texttt{compress} measures $L(H)$
directly via compression ratio; \texttt{mdl} shows the $L(H)$ vs
$L(D \mid H)$ tradeoff in polynomial model selection; \texttt{replicate}
shows what happens when $L(H)$ is underreported; \texttt{search} shows
that verification is evaluating $L(D \mid H)$ cheaply; \texttt{update}
shows the Kolmogorov prior doing work data alone cannot; \texttt{autoloop}
runs the full MDL loop and identifies the gap between current AutoResearch
systems and the complete objective.

\subsection{Structure Is Compressibility (\texttt{ar compress})}

Kolmogorov complexity is uncomputable, but lossless compression provides
a computable upper bound. A string that compresses well has low $\K$; one
that does not compress has high $\K$. The compression ratio approximates
$\K(x)/|x|$---the complexity relative to the raw length.

\begin{center}
\begin{tabular}{lrrl}
\toprule
File & Raw bytes & Ratio & Interpretation \\
\midrule
\texttt{repeated\_pattern} & 1009 & 0.036 & highly structured \\
\texttt{genome\_like}      &  766 & 0.123 & highly structured \\
\texttt{source\_code}      &  868 & 0.414 & moderate structure \\
\texttt{english\_prose}    & 1501 & 0.502 & moderate structure \\
\texttt{pi\_digits}        & 1009 & 0.520 & moderate structure \\
\texttt{random\_noise}     &  369 & 0.564 & low structure \\
\texttt{natural\_law}      &  591 & 0.604 & low structure \\
\bottomrule
\end{tabular}
\end{center}

The genome compresses to 12.3\% of its raw size. DNA looks complex but
is built from repeated motifs; evolution is a compressive process.

The digits of $\pi$ compress to only 52.0\%. $\pi$ has very low true
Kolmogorov complexity---there exists a short program that generates it
exactly---but gzip finds only local repetitions. This illustrates the
gap between the compression proxy and the true $\K(x)$: gzip is an
upper bound, not the true complexity.

Natural laws compress to 60.4\%, below random noise. The structure
in natural laws is semantic, not syntactic; gzip cannot see it. The
true $\K$ of Newton's law---the shortest program predicting all
observations it covers---is extremely low.

\subsection{The MDL Tradeoff Is Real (\texttt{ar mdl})}

MDL model selection fits polynomials of degree 1 through 8 to data
from a known degree, computing $L(H) + L(D \mid H)$ for each.

On a true degree-2 parabola, 40 points, $\sigma = 0.3$:

\begin{center}
\begin{tabular}{rrrr}
\toprule
Degree & $L(H)$ & $L(D \mid H)$ & Total \\
\midrule
1 & 64.0 & 101.2 & 165.2 \\
2 & 96.0 & $-0.2$ & \textbf{95.8} \\
3 & 128.0 & $-0.4$ & 127.6 \\
4 & 160.0 & $-0.4$ & 159.6 \\
\bottomrule
\end{tabular}
\end{center}

MDL recovers the true degree. The decisive test is
\texttt{overfit\_trap}: same generator, only 12 noisy points
($\sigma = 1.2$):

\begin{center}
\begin{tabular}{rrrr}
\toprule
Degree & $L(H)$ & $L(D \mid H)$ & Total \\
\midrule
1 & 64.0 & 31.4 & \textbf{95.4} \\
2 & 96.0 & 25.4 & 121.4 \\
3 & 128.0 & 24.9 & 152.9 \\
\bottomrule
\end{tabular}
\end{center}

MDL chooses degree 1, not the true degree 2. This is not a failure.
With 12 noisy points, the improvement in fit does not justify the
32-bit complexity penalty. The framework is being epistemically honest:
we do not have enough data to establish the true degree. Science should
say the same thing.

\subsection{The Replication Crisis Is One Missing Term
  (\texttt{ar replicate})}

Simulating p-hacking on null data:

\begin{center}
\begin{tabular}{rrrrl}
\toprule
Tests $N$ & Evidence (bits) & Search cost $\log_2 N$ & Corrected & Verdict \\
\midrule
10  & 5.13 & 3.32 & 1.80    & not significant \\
50  & 7.28 & 5.64 & 1.63    & not significant \\
100 & 6.28 & 6.64 & $-0.37$ & not significant \\
\bottomrule
\end{tabular}
\end{center}

Every naive result clears the 4.32-bit threshold ($p < 0.05$). Every
corrected result is noise. A single test on data with a real effect of
size 0.8 produces 16.45 bits; at $N = 100$ tests the corrected value
is 21.25 bits---the real effect survives with massive margin. The
search cost $\log_2 N$ is exactly the Bonferroni correction in bits.
The replication crisis is the empirical consequence of omitting this
term from $L(H)$.

\subsection{Verification Is Constant; Search Scales (\texttt{ar search})}

\begin{center}
\begin{tabular}{lrr}
\toprule
Candidate & Tests passed & Verify time (ms) \\
\midrule
\texttt{broken\_empty\_return}  & 0/8 & 0.13 \\
\texttt{broken\_mutates\_input} & 0/8 & 0.62 \\
\texttt{broken\_off\_by\_one}   & 0/8 & 0.09 \\
\texttt{correct\_bubble}        & 8/8 & 0.07 \\
\texttt{correct\_recursive}     & 8/8 & 0.08 \\
\bottomrule
\end{tabular}
\end{center}

Verification time is flat. Complexity and correctness do not affect
verification cost. Search cost accumulates at 6.1$\times$ the
verification cost for 5 candidates; with 50,000 candidates the ratio
is ${\sim}60{,}000\times$. The verification cost never moves. The
candidate \texttt{broken\_returns\_none} is caught in 0.07\,ms with
no knowledge of its source or intent. The object speaks for itself.

\subsection{The Right Prior Converges Faster (\texttt{ar update})}

Four hypotheses---$3x$, $3x + 0{\cdot}x^2$, $3x + 0{\cdot}x^3$,
$3x + 0{\cdot}x^2 + 0{\cdot}x^3$---predict identically for every
input.

\begin{center}
\begin{tabular}{lrr}
\toprule
Step & Uniform on $3x$ & Kolmogorov on $3x$ \\
\midrule
Prior & 25.0\% & 99.6\% \\
1     & 25.0\% & 99.6\% \\
5     & 25.0\% & 99.6\% \\
\bottomrule
\end{tabular}
\end{center}

The uniform prior stays at 25\% forever. The Kolmogorov prior starts
at 99.6\% before any data. Among hypotheses that predict equally well,
the prior decides---and it should be the Kolmogorov prior. In the
weak-signal case (true $2x$, noise $= 1.5$), the Kolmogorov prior
starts $2x$ lower than uniform (17.4\% vs 25\%) yet reaches 78\% at
step 3 versus 68.5\% for uniform. The mechanism: the Kolmogorov prior
suppressed the wrong alternatives from the start.

\subsection{The Karpathy Gap (\texttt{ar autoloop})}

The \texttt{autoloop} experiment models the Karpathy AutoResearch loop
in two modes: full MDL ($L(H) + L(D \mid H)$) and $L(D \mid H)$ only.
Starting from a baseline that already passes all 8 tests:

\begin{center}
\begin{tabular}{llrrrrl}
\toprule
Iter & Mutation & Tests & $L(H)$ & $L(D|H)$ & MDL & Decision \\
\midrule
0 & baseline                      & 8/8 & 279 &  0.0 & 279.0 & --- \\
1 & \texttt{01\_fix\_break}        & 8/8 & 257 &  0.0 & 257.0 & keep \\
2 & \texttt{02\_add\_type\_check}  & 8/8 & 315 &  0.0 & 315.0 & discard \\
3 & \texttt{03\_simplify}          & 8/8 & 170 &  0.0 & 170.0 & keep \\
4 & \texttt{04\_introduce\_bug}    & 3/8 & 197 & 15.0 & 212.0 & discard \\
5 & \texttt{05\_use\_builtin}      & 8/8 &  97 &  0.0 &  97.0 & keep \\
6 & \texttt{06\_overengineered}    & 8/8 & 211 &  0.0 & 211.0 & discard \\
\bottomrule
\end{tabular}
\end{center}

Full MDL: baseline $279.0 \to 97.0$, \textbf{65.2\% reduction}.
$L(D \mid H)$-only: 0 improvements, stuck forever. The decisive case
is \texttt{06\_overengineered}: passes all tests, $L(D \mid H) = 0$,
$L(H) = 211$. Full MDL discards it; $L(D \mid H)$-only cannot
distinguish it from \texttt{05\_use\_builtin}.

This is the Karpathy gap: Karpathy's \texttt{val\_bpb} is
$L(D \mid H)$ in QBBN coding (Section~\ref{sec:coding}) without
$L(H)$. The gap is not a toy artifact. Coppola~\cite{coppola2026}
proves that sigmoid transformers implement exact BP on
QBBN-structured factor graphs, so every transformer trained to minimize
validation loss without an explicit complexity penalty is running the
same incomplete objective. The missing term is $L(H)$. Adding it
completes MDL and prevents accumulating complexity over hundreds of
iterations whenever doing so does not hurt validation loss.
\section{Experiments II: Logical Propositions}
\label{sec:experiments-logic}

The experiments in this section instantiate the MDL objective function
in the domain of logical proposition datasets, using the QBBN coding
scheme developed in Section~\ref{sec:coding}. All results are
reproducible from the \texttt{mdl-science} repository.

A dataset is a list of boolean proposition observations. A theory is a
set of \emph{implication links} (premise $\to$ conclusion, firing when
the premise is true) and \emph{conjunctive links}
($[\text{premise}_1, \ldots, \text{premise}_n] \to$ conclusion, firing
only when all premises are true). These correspond directly to the QBBN:
conjunctive links are AND gates, implication links are OR-gate inputs.

The two-part code is concrete. $L(H)$ counts the number of links times
the bits per weight---a fixed cost per link, independent of dataset
size. $L(D \mid H)$ is the total negative log-likelihood of the observed
propositions under the theory's conditional probabilities, with link
weights calibrated empirically from the data. Propositions not covered
by any link fall back to their empirical base rate, ensuring the theory
can only help---never hurt---on uncovered propositions.

\subsection{Conjunctive Structure Wins (\texttt{mdl compress})}

Data where $r = p \wedge q$ exactly ($n = 50$ triples). Six theories
compete. Results are reported as MDL score and bits saved versus null
(positive = beats no theory; negative = worse than no theory):

\begin{center}
\begin{tabular}{lrrrr}
\toprule
Theory & $L(H)$ & $L(D \mid H)$ & MDL & Bits saved \\
\midrule
null (no links)                 &  0.0 & 143.5 & 143.5 &      $0.0$ \\
one link: $p \to r$             &  8.0 & 134.2 & 142.2 &    $+1.3$ \\
one link: $q \to r$             &  8.0 & 140.0 & 148.0 &    $-4.5$ \\
two links: $p \to r$, $q \to r$ & 16.0 & 141.8 & 157.8 &   $-14.3$ \\
conjunctive: $p \wedge q \to r$ &  8.0 & 117.1 & 125.1 & $+\mathbf{18.4}$ \\
wrong link: $p \to q$           &  8.0 & 143.2 & 151.2 &    $-7.7$ \\
\bottomrule
\end{tabular}
\end{center}

\textit{Bits saved = MDL(null) $-$ MDL(theory).}

The conjunctive theory wins by 18.4 bits using only 8 bits to describe
one link. The disjunctive two-link theory loses by 14.3 bits despite
having both premises available: it misfires when only one of $p$ or $q$
is true, predicting $r = 1$ when $r = 0$. The wrong link loses by 7.7
bits. Even the partial theory $p \to r$ barely wins at $+1.3$ bits.

The key result: having more links is not always better. The two-link
disjunctive theory has strictly more information about $r$ than the
one-link conjunctive theory, yet it performs worse. What matters is
whether the links capture the true generative structure. The disjunctive
theory has the right variables but the wrong logical connective.

This result is a direct consequence of the QBBN coding scheme. The AND
gate (conjunctive link) is a zero-parameter, deterministic structure
that adds only its own description to $L(H)$. The OR-gate inputs
(implication links) each add weight parameters. When the true structure
is conjunctive, the conjunctive coding is maximally efficient: it pays
zero for the AND computation itself, only for the single weight that
says how strongly the conjunction predicts the conclusion.

\subsection{Theories Are Amortized Over Data (\texttt{mdl scale})}

One conjunctive link, $L(H) = 8$ bits, regularity $= 1.0$, $n$ from
5 to 1000:

\begin{center}
\begin{tabular}{rrrrr}
\toprule
$n$ & $L(D)$ & MDL & Bits saved & Comp.\% \\
\midrule
5    &   13.3 &   18.7 &   $-5.4$ & $-40.2$\% \\
10   &   28.8 &   31.6 &   $-2.8$ &  $-9.7$\% \\
25   &   70.9 &   65.8 &   $+5.0$ &  $+7.1$\% \\
50   &  136.7 &  118.1 &  $+18.5$ & $+13.6$\% \\
100  &  276.2 &  231.2 &  $+44.9$ & $+16.3$\% \\
200  &  535.7 &  440.5 &  $+95.2$ & $+17.8$\% \\
500  & 1326.8 & 1079.0 & $+247.9$ & $+18.7$\% \\
1000 & 2680.7 & 2170.1 & $+510.6$ & $+19.0$\% \\
\bottomrule
\end{tabular}
\end{center}

\textit{Bits saved = $L(D) -$ MDL. Positive means the theory beats the
no-theory baseline.}

$L(H) = 8$ bits at every scale. The theory loses at $n = 5$ and
$n = 10$, crosses into positive territory between $n = 10$ and $n = 25$,
and saves 510 bits at $n = 1000$---64 times its own description length.

This is the formal version of the intuition behind scientific laws.
Newton's laws cost a fixed number of bits to state; every new
observation is compressed for free once the theory exists. The
amortization is not a metaphor. It is a computable quantity,
demonstrated here at the level of individual bits, with the crossover
point (theory overhead recovered at $\approx n = 15$ pairs) precisely
measurable.

The QBBN coding scheme makes this particularly clean. Because $L(H)$
counts only the structural description of the links (which propositions,
which logical connective, what weight precision), it is genuinely
independent of dataset size. The same 8 bits describes the theory
whether the dataset has 5 observations or 5 million. This is the coding
scheme doing its job: separating the cost of the theory from the cost
of the data.

\subsection{MDL Selects Correct Theory Complexity (\texttt{mdl compete})}

Four theories at three regularity levels. T0 = null, T1 = one
implication link, T2 = correct conjunctive link, T3 = T2 plus one
spurious implication link:

\begin{center}
\begin{tabular}{lrrrr}
\toprule
Theory & $L(H)$ & $L(D \mid H)$ & MDL & Bits saved \\
\midrule
\multicolumn{5}{l}{\textit{regularity = 1.0}} \\
T0: null                &  0.0 & 276.2 & 276.2 &       $0.0$ \\
T1: one link            &  8.0 & 272.6 & 280.6 &      $-4.5$ \\
T2: correct conjunctive &  8.0 & 223.2 & 231.2 & $+\mathbf{44.9}$ \\
T3: T2 + spurious link  & 16.0 & 223.2 & 239.2 &     $+36.9$ \\
\midrule
\multicolumn{5}{l}{\textit{regularity = 0.5}} \\
T0: null                &  0.0 & 272.5 & 272.5 & $\mathbf{0.0}$ \\
T1: one link            &  8.0 & 272.0 & 280.0 &      $-7.6$ \\
T2: correct conjunctive &  8.0 & 262.6 & 270.6 &      $+1.9$ \\
T3: T2 + spurious link  & 16.0 & 262.6 & 278.6 &      $-6.1$ \\
\midrule
\multicolumn{5}{l}{\textit{regularity = 0.0}} \\
T0: null                &  0.0 & 273.5 & 273.5 & $\mathbf{0.0}$ \\
T1: one link            &  8.0 & 273.5 & 281.5 &      $-8.0$ \\
T2: correct conjunctive &  8.0 & 273.5 & 281.5 &      $-8.0$ \\
T3: T2 + spurious link  & 16.0 & 273.5 & 289.5 &     $-16.0$ \\
\bottomrule
\end{tabular}
\end{center}

\textit{Bits saved = MDL(T0) $-$ MDL(theory).}

At regularity $= 1.0$, T2 saves 44.9 bits. T3 saves 36.9 bits---
exactly 8 bits less than T2. The spurious link costs exactly its
description length and changes $L(D \mid H)$ not at all. This is
Occam's razor in its precise form: \emph{a theory that is correct but
unnecessarily complex loses to the minimal correct theory by exactly
the description length of the excess.}

At regularity $= 0.5$, T2 barely saves 1.9 bits. The signal is real
but weak; T3 cannot recover the extra 8-bit cost. At regularity $=
0.0$, T0 wins; T1 and T2 tie at exactly $-8.0$ bits; T3 pays $-16.0$
bits. The framework penalizes each spurious link by exactly its
description length.

The regularity crossover is the MDL phase transition of
Section~\ref{sec:connections} instantiated at the bit level. As
regularity increases from 0 to 1, the optimal theory switches
discontinuously from T0 to T2. The switch is not gradual: T2 goes
from $-8.0$ bits (noise) to $+1.9$ bits (weak signal) to $+44.9$ bits
(perfect regularity). The optimal hypothesis class jumps when the
theory's amortized savings first exceed its description length cost.
This is exactly what Kuhn observed in the history of science and what
Lakatos formalized as the distinction between progressive and
degenerating research programmes. Here it is a measurable threshold
in bits.
\section{Verification and the Present Moment}

\subsection{The Search/Verify Asymmetry}

The objective function tells us what science is optimizing. Running the
optimization loop requires solving a different problem: given a proposed
hypothesis $H$, how do you verify that it actually advances the objective?

The key structural fact:
\begin{quote}
  Finding a good $H$ is hard. Verifying a proposed $H$ is cheap.
\end{quote}
This asymmetry is not incidental. It is a feature of the Kolmogorov
framework, and it is the design principle behind every sound verification
protocol. Finding the shortest program for $x$ requires searching an infinite
space. Checking whether a proposed program outputs $x$ requires only running
it.

\subsection{Algorithmic Verification}

For claims in the verifiable domain, verification can be automated.

\textbf{Formal proof checking.} If a hypothesis is expressed as a theorem
with a machine-checkable proof (Lean~4, Coq, Isabelle), verification is
linear in proof length. The checker evaluates the \emph{object}, not the
source. A proof submitted by an unknown agent is as valid as one submitted
by a Fields Medalist, if it type-checks. This is the purest instantiation of
the Kolmogorov principle.

\textbf{Empirical reproduction.} Given code, weights, and data, rerun the
experiment. Does the reported metric reproduce? This is cheap relative to
the original research cost and provides a binary signal.

\textbf{Compression-based verification.} Does the proposed theory compress
the data better than the incumbent? MDL makes this a computable number.
Given a fixed coding scheme, compute $L(H) + L(D \mid H)$ for the new theory
and for the incumbent. If the new value is lower, the theory advances the
objective. No subjective judgment required.

\subsection{The Untrusted Source Principle}

The most important design principle for any verification system:
\begin{quote}
  The verifier does not need to trust the source.
\end{quote}
Under the Kolmogorov framework, a hypothesis is evaluated on its length and
its predictive accuracy---properties of the object itself. The identity,
credentials, or intentions of whoever proposed it are irrelevant to its truth.

This is an engineering requirement, not just a philosophical point. A
distributed AutoResearch system cannot verify the identity or intent of every
contributing agent. But if verification is a property of the object (does the
proof type-check? does the loss decrease? does the description length shrink?),
the system can accept submissions from arbitrary untrusted sources and remain
sound. The math is the trust anchor.

\subsection{Why This Matters Now}

Two developments make the framework more than a philosophical exercise.

\textbf{Autonomous AI research.} Systems like Karpathy's AutoResearch close
the optimization loop at the level of ML experiments: propose a modification,
verify whether it reduces the loss, keep it if it does, repeat. This is
Solomonoff induction approximated by the only means available: a computable,
distributed, heuristic search over hypothesis space, with cheap verification
at each step.

\textbf{Machine-verified formal proof.} Large proofs in Lean~4 and Coq are
being produced and verified at scale. The verification is automated; the
search (finding the proof) is expensive; the check (type-checking) is cheap.
The asymmetry is the same.

Both represent the first time the optimization loop of science can be
partially automated. The objective function---minimize $L(H) + L(D \mid H)$---
is the same as it has always been. What is new is the infrastructure to
approximate it at scale.

\subsection{Conclusion}

The history of philosophy of science is a record of humanity circling a
single structure---the problem of optimal inference in a computable
world---without the mathematical tools to see it directly. Hume identified
the core impossibility. Popper, Kuhn, and Lakatos each described a facet
of the solution without being able to state it. The Bayesian tradition
formalized the update rule. Kolmogorov and Solomonoff formalized the prior.

The tools now exist to state the full solution. Science is the distributed
search for hypotheses that minimize description length. Verification is cheap
relative to discovery. The prior is the Solomonoff prior. The update rule is
Bayes' rule. The objective function is MDL.

Everything else is details.
\section{Automated Theorem Proving Is Science}
\label{sec:agi}

\subsection{The Critics' Claim}

A recurring objection in discussions of artificial general intelligence runs
as follows. Automated theorem proving---even at superhuman scale, even
across all of mathematics---is not AGI, because formal reasoning is not the
same as doing science. Science requires engaging with the empirical world:
forming hypotheses, running experiments, updating on data. A system that
only manipulates symbols inside a formal system is doing something
categorically different from what scientists do. Theorem proving, however
impressive, is a party trick. It is not intelligence.

This objection is widespread and sounds compelling. We argue it is wrong,
and that it is wrong for a precise reason: it rests on an implicit theory of
science that has never been written down---and when you write it down, it
does not survive scrutiny.

\subsection{The Implicit Theory Being Assumed}

The critics' argument requires a theory of what science is, specifically one
that excludes formal reasoning. The theory being assumed, roughly, is this:
science is the activity of interacting with the physical world through
observation and experiment, and formal reasoning is merely a tool that
science uses, not science itself.

This theory has intuitive appeal but no formal content. It cannot tell you:
\begin{itemize}
  \item When two competing theories are equally consistent with the data,
        which is better?
  \item How much complexity is a theory allowed before it is no longer
        parsimonious?
  \item When should a paradigm be abandoned in favor of a competitor?
  \item What counts as a legitimate auxiliary hypothesis versus an ad hoc
        patch?
\end{itemize}
These are the central questions of the philosophy of science, and the
informal ``interaction with the physical world'' theory has no answers to
them. It is not a theory of science. It is a description of one setting in
which science sometimes happens.

\subsection{The Formal Theory}

The Bayesian/Kolmogorov framework provides an actual theory of science,
with formal answers to all of the above questions. Restating it precisely:

\begin{definition}[Scientific activity]
  A process is engaged in scientific activity if it:
  \begin{enumerate}
    \item Maintains a space of computable hypotheses $\mathcal{H}$;
    \item Assigns prior probabilities $P(H) \propto 2^{-\K(H)}$ to
          hypotheses, preferring shorter descriptions;
    \item Evaluates hypotheses against observations via a likelihood
          function $P(D \mid H)$;
    \item Updates beliefs by Bayes' rule: $P(H \mid D) \propto
          P(D \mid H) \cdot P(H)$; and
    \item Selects hypotheses that minimize $L(H) + L(D \mid H)$.
  \end{enumerate}
\end{definition}

Notice what is \emph{not} in this definition: laboratories, physical
instruments, empirical data in any particular form, or indeed any
requirement that $D$ come from the external world rather than from a
formal system. The definition is about the \emph{structure} of the
inference, not the \emph{substrate} in which it occurs.

\subsection{The Verification Oracle Is Interchangeable}

The key insight is that the search/verify asymmetry
(Section~\ref{sec:verification}) holds regardless of what the verifier is.

In empirical science, the verifier is experiment: you propose a hypothesis,
derive a prediction, run the experiment, and check whether the prediction
holds. Finding a good hypothesis is hard; running the experiment is cheap
relative to the search.

In formal mathematics, the verifier is the proof checker: you propose a
theorem and a proof, and the type-checker verifies it in time linear in
proof length. Finding a good proof is hard (it may require deep insight or
exhaustive search); checking it is cheap.

Both are instances of the same abstract structure:
\[
  \text{Search}(H \in \mathcal{H}) \xrightarrow{\text{expensive}}
  \hat{H} \xrightarrow{\text{cheap verification}} \text{accept or reject}
\]
The oracle that implements ``cheap verification'' differs in its physical
instantiation---a particle accelerator versus a type-checker---but the
computational structure is identical. A system that can search hypothesis
space and verify proposals is doing science. The nature of the oracle is
irrelevant to this classification.

\subsection{What Formal Mathematics Looks Like Under This Theory}

Under the Kolmogorov framework, a mathematical theorem is a hypothesis about
the structure of a formal system. A proof is the verification. The
``observations'' $D$ are the axioms and previously established theorems. The
``experiment'' is type-checking.

This is not a strained analogy. Consider what a mathematician does:
\begin{itemize}
  \item She maintains a space of possible theorems (hypotheses).
  \item She assigns higher prior probability to theorems that are short,
        elegant, and follow from a few strong principles---precisely
        $P(H) \propto 2^{-\K(H)}$.
  \item She evaluates a candidate theorem against the existing body of
        mathematics (the ``data'').
  \item She updates: a theorem that follows from deep principles and
        connects many areas is more valuable than one that requires a
        long, ad hoc proof.
  \item She publishes results that compress the existing body of knowledge
        into shorter descriptions---unifying previously separate results,
        revealing that two seemingly different theorems are instances of
        one.
\end{itemize}
This is the MDL objective, applied to a formal domain. The mathematician
is doing exactly what the physicist does. The substrate differs; the
objective function is the same.

\subsection{The Objection Answered}

The critics are right that theorem proving, as typically practiced, is not
the same as doing science. But this is because automated theorem provers as
currently deployed do not have the right objective function. They are
optimized to prove specified theorems, not to search for theorems worth
proving. They are given $H$ and asked to verify it, rather than searching
$\mathcal{H}$ for the $H$ that minimizes description length.

This is an engineering gap, not a conceptual impossibility. A theorem-proving
system equipped with:
\begin{enumerate}
  \item A prior over theorems weighted by description length,
  \item A likelihood that measures how much a theorem compresses and
        unifies the existing mathematical corpus,
  \item A search procedure over the space of provable statements, and
  \item A cheap verifier (the type-checker),
\end{enumerate}
would be doing science in the full formal sense. It would propose conjectures
that compress the mathematical corpus, verify them, keep the ones that
advance the MDL objective, and discard the rest. This is Solomonoff
induction instantiated in a formal domain.

Such a system would not be doing something categorically different from what
Einstein did when he proposed general relativity, or what Ramanujan did when
he produced his notebooks. It would be doing the same thing, faster and
without biological constraints.

\subsection{AGI and the Objective Function}

The question of whether a system constitutes artificial general intelligence
is often posed in terms of capabilities: can it reason, plan, learn, adapt?
We propose a different framing, grounded in the framework of this paper.

\begin{proposition}[AGI as correct objective function]
  A system exhibits artificial general intelligence if and only if it
  implements---exactly or approximately---Solomonoff induction: maintaining
  a prior over computable hypotheses weighted by description length and
  updating by Bayes' rule on available observations.
\end{proposition}

This definition has several virtues. It is formal and precise. It is
substrate-independent: the system can operate on physical observations,
mathematical structures, or any other computable domain. It is graded:
systems can approximate Solomonoff induction more or less closely, with
better approximations constituting more general intelligence. And it
connects AGI directly to the objective function for science: a system that
is doing science in the formal sense is, by this definition, exhibiting
general intelligence.

Under this definition, the question ``is theorem proving AGI?'' has a
precise answer: it depends on whether the theorem-proving system has the
right objective function. A system that is given theorems to prove and
proves them is not AGI---it is a very capable special-purpose tool. A
system that searches for theorems worth proving, verifies them, and updates
its prior based on what it finds---minimizing description length over the
space of mathematical knowledge---is AGI, regardless of whether it ever
runs a physical experiment.

\subsection{The Broader Implication}

The critics who dismiss theorem proving as ``not really AGI'' are, without
knowing it, making an implicit claim about the objective function for
science. They are claiming that the objective function requires physical
observation as input. We have argued that this claim is false: the objective
function is MDL, the input can be any observations from any verifiable
domain, and the structure of inference is the same regardless.

The tools to build such a system exist. The objective function is written
down---in this paper and in the century of work it synthesizes. The
verification infrastructure exists: Lean~4 type-checks proofs in linear
time; benchmark evaluation checks model performance in seconds; experimental
platforms check empirical predictions in hours or days. The search
infrastructure is being built: AutoResearch~\cite{karpathy2026},
agent-based research systems, and large language models used as hypothesis
generators are all approximating the search step.

What has been missing is the theoretical unification: the recognition that
all of these are the same process, that they all approximate the same
objective function, and that a system implementing that objective function
in any verifiable domain is doing science and exhibiting general
intelligence.

That recognition is the contribution of this paper.

\bibliographystyle{plain}
\bibliography{refs}

\end{document}